% GENERATED FILE. Edit canonical lessons, then run npm run generate:books.
% canonical-source-hash: fnv1a64:2a9f3de1cf0d2d0c
% canonical-lessons: KA-C424-avalokisu, KA-C424-vimarsisu, KA-C424-suttole, KA-C424-nirdista, KA-C424-parisodhane

\chapter{Checks and Reviews}
\label{ch:ka-checks-and-reviews}
\hlchaptermodality{424}

\begin{chapteropening}
\textbf{By the end of this chapter:} \emph{I can say to observe, to critique, a circular, specific and an audit.}

Complete the last lesson of chapter 424: I can say to observe, to critique, a circular, specific and an audit.
\end{chapteropening}

\section[\texorpdfstring{\kn{ಅವಲೋಕಿಸು} (avalōkisu)}{ಅವಲೋಕಿಸು (avalōkisu)}]{\kn{ಅವಲೋಕಿಸು} (avalōkisu) — to observe; to review}

\label{lesson:KA-C424-avalokisu}



\begin{quote}
Before the new one: say the Kannada for a trick, then the Kannada for to divide.
\end{quote}

\subsection*{You'll want to know: \kn{ಅವಲೋಕಿಸು}}
\textbf{\kn{ಅವಲೋಕಿಸು}} — \emph{avalōkisu} — \textquotedblleft{}to observe; to review\textquotedblright{}.

\begin{quote}
\kn{ವರ್ಷದ} \kn{ಕೊನೆಯಲ್ಲಿ} \kn{ನಾನು} \kn{ನನ್ನ} \kn{ಖರ್ಚನ್ನು} \kn{ಅವಲೋಕಿಸಿದೆ}.
\end{quote}

\begin{quote}
\emph{varṣada koneyalli nānu nanna kharcannu avalōkiside.}
\end{quote}

\begin{quote}
At the end of the year I reviewed my spending.
\end{quote}

Past (I …): \kn{ಅವಲೋಕಿಸಿದೆ}, \emph{avalōkiside}.

\kn{ಅವಲೋಕನ}, a review, an overview.

\begin{cousinweb}
\textbf{\kn{ಅವಲೋಕಿಸು}} (\emph{avalōkisu}) is built on Sanskrit \emph{avalokana}, \textquotedblleft{}seeing, viewing, observing\textquotedblright{} (Monier-Williams), with the Kannada verb ending -\emph{isu}.
\end{cousinweb}

\begin{grammarlens}[title={What you've built}]
One more word for saying what happens if.
\end{grammarlens}

\subsection*{Your turn}
\begin{itemize}
\raggedright
  \item \emph{Say it:} \emph{avalōkisu}
  \item \emph{Say it:} \emph{avalōkisu}, once more
  \item \emph{Say it:} \emph{varṣada koneyalli nānu nanna kharcannu avalōkiside.} — At the end of the year I reviewed my spending.
  \item \emph{Recall:} say the Kannada for a trick, then the Kannada for to divide, then say \emph{avalōkisu} again
\end{itemize}

\begin{tcolorbox}[breakable,colback=teal!4,colframe=teal!35!black,title={Before you move on}]
What does \kn{ಅವಲೋಕಿಸು} mean? (\textquotedblleft{}To observe; to review\textquotedblright{}.) Say it once more.
\end{tcolorbox}
\section[\texorpdfstring{\kn{ವಿಮರ್ಶಿಸು} (vimarśisu)}{ವಿಮರ್ಶಿಸು (vimarśisu)}]{\kn{ವಿಮರ್ಶಿಸು} (vimarśisu) — to critique}

\label{lesson:KA-C424-vimarsisu}



\begin{quote}
Before the new one: say the Kannada for to divide, then the Kannada for to observe.
\end{quote}

\subsection*{You'll want to know: \kn{ವಿಮರ್ಶಿಸು}}
\textbf{\kn{ವಿಮರ್ಶಿಸು}} — \emph{vimarśisu} — \textquotedblleft{}to critique\textquotedblright{}.

\begin{quote}
\kn{ಪತ್ರಿಕೆ} \kn{ಹೊಸ} \kn{ಪುಸ್ತಕವನ್ನು} \kn{ವಿಮರ್ಶಿಸಿತು}.
\end{quote}

\begin{quote}
\emph{patrike hosa pustakavannu vimarśisitu.}
\end{quote}

\begin{quote}
The newspaper reviewed the new book.
\end{quote}

Past (I …): \kn{ವಿಮರ್ಶಿಸಿದೆ}, \emph{vimarśiside}.

\kn{ವಿಮರ್ಶೆ}, a review; \kn{ವಿಮರ್ಶಕ}, a critic.

\begin{cousinweb}
\textbf{\kn{ವಿಮರ್ಶಿಸು}} (\emph{vimarśisu}) is built on Sanskrit \emph{vimarśa}, \textquotedblleft{}consideration, critical test\textquotedblright{} (Monier-Williams), with the Kannada verb ending -\emph{isu}.
\end{cousinweb}

\begin{grammarlens}[title={What you've built}]
One more word for saying what happens if.
\end{grammarlens}

\subsection*{Your turn}
\begin{itemize}
\raggedright
  \item \emph{Say it:} \emph{vimarśisu}
  \item \emph{Say it:} \emph{vimarśisu}, once more
  \item \emph{Say it:} \emph{patrike hosa pustakavannu vimarśisitu.} — The newspaper reviewed the new book.
  \item \emph{Recall:} say the Kannada for to divide, then the Kannada for to observe, then say \emph{vimarśisu} again
\end{itemize}

\begin{tcolorbox}[breakable,colback=teal!4,colframe=teal!35!black,title={Before you move on}]
What does \kn{ವಿಮರ್ಶಿಸು} mean? (\textquotedblleft{}To critique\textquotedblright{}.) Say it once more.
\end{tcolorbox}
\section[\texorpdfstring{\kn{ಸುತ್ತೋಲೆ} (suttōle)}{ಸುತ್ತೋಲೆ (suttōle)}]{\kn{ಸುತ್ತೋಲೆ} (suttōle) — a circular}

\label{lesson:KA-C424-suttole}



\begin{quote}
Before the new one: say the Kannada for to observe, then the Kannada for to critique.
\end{quote}

\subsection*{You'll want to know: \kn{ಸುತ್ತೋಲೆ}}
\textbf{\kn{ಸುತ್ತೋಲೆ}} — \emph{suttōle} — \textquotedblleft{}a circular\textquotedblright{}.

\begin{quote}
\kn{ರಜೆಯ} \kn{ಬಗ್ಗೆ} \kn{ಶಾಲೆ} \kn{ಒಂದು} \kn{ಸುತ್ತೋಲೆ} \kn{ಕಳುಹಿಸಿತು}.
\end{quote}

\begin{quote}
\emph{rajeya bagge śāle ondu suttōle kaḷuhisitu.}
\end{quote}

\begin{quote}
The school sent a circular about the holiday.
\end{quote}

\kn{ಸುತ್ತು} (to go round) + \kn{ಓಲೆ} (a letter).

\begin{cousinweb}
\textbf{\kn{ಸುತ್ತೋಲೆ}} (\emph{suttōle}) is native Dravidian, built from \emph{suttu}, \textquotedblleft{}to go round\textquotedblright{} (DEDR 2715), and \emph{ōle}, \textquotedblleft{}a palm leaf written on, a letter\textquotedblright{} (DEDR 1070).
\end{cousinweb}

\begin{grammarlens}[title={What you've built}]
One more word for saying what happens if.
\end{grammarlens}

\subsection*{Your turn}
\begin{itemize}
\raggedright
  \item \emph{Say it:} \emph{suttōle}
  \item \emph{Say it:} \emph{suttōle}, once more
  \item \emph{Say it:} \emph{rajeya bagge śāle ondu suttōle kaḷuhisitu.} — The school sent a circular about the holiday.
  \item \emph{Recall:} say the Kannada for to observe, then the Kannada for to critique, then say \emph{suttōle} again
\end{itemize}

\begin{tcolorbox}[breakable,colback=teal!4,colframe=teal!35!black,title={Before you move on}]
What does \kn{ಸುತ್ತೋಲೆ} mean? (\textquotedblleft{}A circular\textquotedblright{}.) Say it once more.
\end{tcolorbox}
\section[\texorpdfstring{\kn{ನಿರ್ದಿಷ್ಟ} (nirdiṣṭa)}{ನಿರ್ದಿಷ್ಟ (nirdiṣṭa)}]{\kn{ನಿರ್ದಿಷ್ಟ} (nirdiṣṭa) — specific; definite}

\label{lesson:KA-C424-nirdista}



\begin{quote}
Before the new one: say the Kannada for to critique, then the Kannada for a circular.
\end{quote}

\subsection*{You'll want to know: \kn{ನಿರ್ದಿಷ್ಟ}}
\textbf{\kn{ನಿರ್ದಿಷ್ಟ}} — \emph{nirdiṣṭa} — \textquotedblleft{}specific; definite\textquotedblright{}.

\begin{quote}
\kn{ಅವನಿಗೆ} \kn{ಯಾವುದೇ} \kn{ನಿರ್ದಿಷ್ಟ} \kn{ಗುರಿ} \kn{ಇಲ್ಲ}.
\end{quote}

\begin{quote}
\emph{avanige yāvudē nirdiṣṭa guri illa.}
\end{quote}

\begin{quote}
He has no definite goal.
\end{quote}

\kn{ನಿರ್ದಿಷ್ಟವಾಗಿ}, specifically.

\begin{cousinweb}
\textbf{\kn{ನಿರ್ದಿಷ್ಟ}} (\emph{nirdiṣṭa}) is borrowed from Sanskrit \emph{nirdiṣṭa}, \textquotedblleft{}pointed out, indicated\textquotedblright{} (Monier-Williams).
\end{cousinweb}

\begin{grammarlens}[title={What you've built}]
One more word for saying what happens if.
\end{grammarlens}

\subsection*{Your turn}
\begin{itemize}
\raggedright
  \item \emph{Say it:} \emph{nirdiṣṭa}
  \item \emph{Say it:} \emph{nirdiṣṭa}, once more
  \item \emph{Say it:} \emph{avanige yāvudē nirdiṣṭa guri illa.} — He has no definite goal.
  \item \emph{Recall:} say the Kannada for to critique, then the Kannada for a circular, then say \emph{nirdiṣṭa} again
\end{itemize}

\begin{tcolorbox}[breakable,colback=teal!4,colframe=teal!35!black,title={Before you move on}]
What does \kn{ನಿರ್ದಿಷ್ಟ} mean? (\textquotedblleft{}Specific; definite\textquotedblright{}.) Say it once more.
\end{tcolorbox}
\section[\texorpdfstring{\kn{ಪರಿಶೋಧನೆ} (pariśōdhane)}{ಪರಿಶೋಧನೆ (pariśōdhane)}]{\kn{ಪರಿಶೋಧನೆ} (pariśōdhane) — an audit; a scrutiny}

\label{lesson:KA-C424-parisodhane}



\begin{quote}
Before the new one: say the Kannada for a circular, then the Kannada for specific.
\end{quote}

\subsection*{You'll want to know: \kn{ಪರಿಶೋಧನೆ}}
\textbf{\kn{ಪರಿಶೋಧನೆ}} — \emph{pariśōdhane} — \textquotedblleft{}an audit; a scrutiny\textquotedblright{}.

\begin{quote}
\kn{ಕಂಪನಿಯ} \kn{ಲೆಕ್ಕಗಳ} \kn{ಪರಿಶೋಧನೆ} \kn{ನಡೆಯಿತು}.
\end{quote}

\begin{quote}
\emph{kampaniya lekkagaḷa pariśōdhane naḍeyitu.}
\end{quote}

\begin{quote}
The company's accounts were audited.
\end{quote}

\kn{ಪರಿಶೋಧಕ}, an auditor.

\begin{cousinweb}
\textbf{\kn{ಪರಿಶೋಧನೆ}} (\emph{pariśōdhane}) is borrowed from Sanskrit \emph{pariśodhana}, \textquotedblleft{}cleaning, purification; paying off\textquotedblright{} (Monier-Williams).
\end{cousinweb}

\begin{grammarlens}[title={What you've built}]
One more word for saying what happens if.
\end{grammarlens}

\subsection*{Your turn}
\begin{itemize}
\raggedright
  \item \emph{Say it:} \emph{pariśōdhane}
  \item \emph{Say it:} \emph{pariśōdhane}, once more
  \item \emph{Say it:} \emph{kampaniya lekkagaḷa pariśōdhane naḍeyitu.} — The company's accounts were audited.
  \item \emph{Recall:} say the Kannada for a circular, then the Kannada for specific, then say \emph{pariśōdhane} again
\end{itemize}

\begin{tcolorbox}[breakable,colback=teal!4,colframe=teal!35!black,title={Before you move on}]
What does \kn{ಪರಿಶೋಧನೆ} mean? (\textquotedblleft{}An audit; a scrutiny\textquotedblright{}.) Say it once more.
\end{tcolorbox}
